\expandafter\ifx\csname CAPETINCLUDEMODE\endcsname\relax
\documentclass[11pt,a4paper]{article}
% preamble.tex — packages et commandes communes a toutes les fiches CAPET SII IE
\input{glyphtounicode}
\pdfgentounicode=1
\usepackage{cmap}
\usepackage[T1]{fontenc}
\usepackage[utf8]{inputenc}
\usepackage[french]{babel}
% babel[french] active a \begin{document} des raccourcis de ponctuation (:;!?)
% qui cassent la syntaxe to[...] de circuitikz (erreur "+ or - expected") ;
% on les desactive, ce qui ne change que l'espacement fin avant ces signes.
\AtBeginDocument{\shorthandoff{:;!?}}
\usepackage[a4paper, top=1.6cm, bottom=1.6cm, left=1.4cm, right=1.4cm]{geometry}
\setlength{\headheight}{14.5pt}
\usepackage{amsmath, amssymb, amsfonts}
\usepackage{siunitx}
\sisetup{locale=FR, output-decimal-marker={,}, inter-unit-product=\ensuremath{{}\cdot{}},
         group-minimum-digits=4}
% Unites non fournies par siunitx, declarees ici une fois pour toutes (fiches + exercices).
% Les gardes \ifdefined evitent tout conflit avec une redeclaration locale restee dans un
% fichier : ne pas retirer les gardes, ne pas redeclarer ces unites dans les fichiers.
\DeclareSIUnit{\jour}{j}
\ifdefined\var\else\DeclareSIUnit{\var}{var}\fi   % puissance reactive
\ifdefined\vad\else\DeclareSIUnit{\vad}{VAD}\fi   % puissance deformante
\ifdefined\tr\else\DeclareSIUnit{\tr}{tr}\fi      % tour (tr/min)
\usepackage{array, booktabs, multirow, tabularx}
\usepackage{xcolor}
\usepackage[most, breakable]{tcolorbox}
\usepackage{multicol}
\usepackage{enumitem}
\usepackage{mathtools}
\usepackage{lmodern}
% Refonte visuelle 2026 (direction "moderne coloré") : Lato pour le texte,
% maths sans-serif assorties via newtxsf (charge apres amsmath/amssymb).
\usepackage[default]{lato}
\usepackage{newtxsf}
\usepackage{microtype}
\usepackage{fancyhdr}
\usepackage{tikz}
\usetikzlibrary{positioning,arrows.meta}
\usepackage{circuitikz}
\usepackage{lastpage}

% ── Couleurs par domaine (palette modernisée 2026) ────────────────────
\definecolor{colElec}{RGB}{29, 78, 216}        % Électrotechnique — bleu vif
\definecolor{colElecLight}{RGB}{219, 234, 254} % fond bleu clair
\definecolor{colAuto}{RGB}{21, 128, 61}        % Automatique — vert
\definecolor{colMeca}{RGB}{194, 65, 12}        % Mécanique — orange
\definecolor{colTherm}{RGB}{185, 28, 28}       % Thermique — rouge
\definecolor{colQT}{RGB}{232, 244, 233}        % (hérité, plus utilisé par questiontype)
\definecolor{colJury}{RGB}{255, 248, 225}      % fond jury (crème)
\definecolor{colJuryFrame}{RGB}{202, 138, 4}   % bordure/bandeau jury (ambre)
\definecolor{colGrayLight}{RGB}{245, 245, 247} % fond gris neutre
\definecolor{colSeparator}{RGB}{180, 180, 180} % gris séparateur
\definecolor{colExo}{RGB}{8, 145, 178}          % Banque d'exercices — cyan (distinct des fiches)
% ── Encre et gris de la refonte ───────────────────────────────────────
\definecolor{colInk}{RGB}{23, 32, 42}          % texte courant (quasi-noir)
\definecolor{colSlate}{RGB}{83, 98, 118}       % gris ardoise (méta, en-têtes)
\definecolor{colRule}{RGB}{203, 213, 225}      % filets clairs
\definecolor{colJuryBg}{RGB}{255, 248, 225}    % fond crème jury
\definecolor{colPresence}{RGB}{83, 98, 118}    % bloc 5 (présence dans les sujets)
\AtBeginDocument{\color{colInk}}

% ── Style encadré exercice + corrigé (banque d'exercices) ─────────────
\tcbset{
  exoBox/.style={
    breakable,
    enhanced,
    colback=colExo!8!white,
    colframe=colExo,
    colbacktitle=colExo,
    coltitle=white,
    fonttitle=\bfseries\small,
    left=4pt, right=4pt, top=3pt, bottom=3pt,
    before skip=5pt, after skip=5pt,
  }
}

% ── Bascule enonce / corrige (banque d'exercices) ─────────────────────
% \CAPETCORRIGE est defini en ligne de commande (pdflatex "\def\CAPETCORRIGE{1}\input{...}")
% pour compiler la variante corriges d'un meme fichier source ; par defaut (non defini),
% la compilation produit la variante enonces. \CAPETSHOW{enonce}{corrige} choisit la branche.
%
% TROIS MODES (refonte 2026-09-20, manuel de revision) :
%   - aucun drapeau   -> imprime l'enonce            (variante enonces, INCHANGE)
%   - \CAPETCORRIGE   -> imprime le corrige          (variante corriges, INCHANGE)
%   - \CAPETMANUEL    -> imprime l'enonce ET met le corrige de cote pour l'annexe
%
% Le mode manuel stocke chaque corrige dans une macro (\csgdef), sans ecriture
% fichier ni re-tokenisation : le contenu (circuitikz, accents, %) traverse intact.
% L'annexe les rejoue ensuite dans l'ordre via \CAPETANNEXE.
\usepackage{etoolbox}

\newcounter{CAPETexo}          % numero continu d'exercice dans le manuel
\setcounter{CAPETexo}{0}
\newcounter{CAPETannexe}       % compteur de relecture en annexe

% Titre courant de la notion en cours (pose par \CAPETNOTION), sert d'intertitre annexe
\newcommand{\CAPETnotioncode}{}
\newcommand{\CAPETnotiontitre}{}

% \CAPETNOTION{code}{titre} : ouvre une notion dans le manuel (memorise le libelle,
% et marque le rang du premier exercice de la notion pour l'intertitre d'annexe).
\newcommand{\CAPETNOTION}[2]{%
  \renewcommand{\CAPETnotioncode}{#1}%
  \renewcommand{\CAPETnotiontitre}{#2}%
  \csgdef{CAPETnotionstart\the\numexpr\value{CAPETexo}+1\relax}{#1 -- #2}%
}

% ── \degree robuste ───────────────────────────────────────────────────
% Certains fichiers d'exercices utilisent \degree. Ni le preambule ni siunitx
% ne le definissent, et \providecommand local ne suffit pas : sous hyperref,
% \degree existe deja comme commande factice (\ERROR), donc \providecommand
% s'abstient et l'appel echoue a la compilation du manuel (erreur invisible
% en compilation isolee). On le (re)definit donc inconditionnellement ici,
% apres le chargement des packages, pour tout le corpus.
\AtBeginDocument{\providecommand{\degree}{}\renewcommand{\degree}{\ensuremath{^\circ}}}

% Libelle decoratif : les fichiers d'exercices utilisent aussi \CAPETSHOW dans
% leur \fancyhead et leur bloc de titre, pour afficher "(enonces)" ou "(corriges)".
% Ces appels-la ne sont PAS des exercices : en mode manuel ils ne doivent ni
% consommer un numero, ni alimenter l'annexe. On les reconnait a leur premier
% argument, strictement egal au libelle ci-dessous.
\newcommand{\CAPETLIBELLEENONCES}{(\'enonc\'es)}

\ifdefined\CAPETMANUEL
  % --- Mode manuel : enonce en place, corrige differe vers l'annexe ---
  \newcommand{\CAPETSHOW}[2]{%
    \def\CAPETargun{#1}%
    \ifx\CAPETargun\CAPETLIBELLEENONCES
      #1%            appel decoratif : rendu tel quel, hors numerotation
    \else
      \CAPETSHOWEXO{#1}{#2}%
    \fi
  }
  \newcommand{\CAPETSHOWEXO}[2]{%
    \stepcounter{CAPETexo}%
    \csgdef{CAPETcorrige\theCAPETexo}{#2}%
    \edef\CAPETtmpnum{\theCAPETexo}%
    #1%
    \par\vspace{2pt}%
    {\footnotesize\itshape\color{colExo}%
     \hyperref[capet:cor:\CAPETtmpnum]{Corrige de l'exercice \CAPETtmpnum\ $\rightarrow$ p.~\pageref{capet:cor:\CAPETtmpnum}}%
     \label{capet:eno:\CAPETtmpnum}}%
  }
  % \CAPETANNEXE : rejoue tous les corriges captures, dans l'ordre, avec intertitres
  \newcommand{\CAPETANNEXE}{%
    \setcounter{CAPETannexe}{0}%
    \loop\ifnum\value{CAPETannexe}<\value{CAPETexo}%
      \stepcounter{CAPETannexe}%
      \edef\CAPETcurnum{\theCAPETannexe}%
      \ifcsname CAPETnotionstart\CAPETcurnum\endcsname
        \par\medskip
        {\normalsize\bfseries\color{colExo}\csuse{CAPETnotionstart\CAPETcurnum}}%
        \par\nobreak\smallskip
      \fi
      \begin{tcolorbox}[exoBox, breakable,
        title={Corrige de l'exercice \CAPETcurnum\
               \hfill\normalfont\footnotesize
               \hyperref[capet:eno:\CAPETcurnum]{enonce p.~\pageref{capet:eno:\CAPETcurnum}}}]%
      \label{capet:cor:\CAPETcurnum}%
      \csuse{CAPETcorrige\CAPETcurnum}%
      \end{tcolorbox}%
    \repeat
  }
\else
  % --- Modes historiques, comportement strictement inchange ---
  \newcommand{\CAPETSHOW}[2]{\ifdefined\CAPETCORRIGE #2\else #1\fi}
  \newcommand{\CAPETANNEXE}{}
\fi

% ── Difficulté 1 à 3 (étoiles pleines, amssymb) ───────────────────────
\newcommand{\difficulte}[1]{%
  \ifcase#1\relax
  \or $\bigstar$\;niveau 1/3%
  \or $\bigstar\bigstar$\;niveau 2/3%
  \or $\bigstar\bigstar\bigstar$\;niveau 3/3%
  \fi
}

% ── Style encadré extrait de document technique (banque d'exercices) ──
\tcbset{
  docTech/.style={
    breakable,
    enhanced,
    colback=colGrayLight,
    colframe=colSeparator,
    fonttitle=\bfseries\small\itshape,
    coltitle=black,
    colbacktitle=colGrayLight,
    left=4pt, right=4pt, top=3pt, bottom=3pt,
    before skip=4pt, after skip=4pt,
  }
}

% ── Style tcolorbox de section principale ────────────────────────────
% Refonte 2026 : bandeau titre coloré arrondi + fond très pâle, sans cadre.
\tcbset{
  ficheSection/.style={
    breakable, enhanced jigsaw,
    frame hidden, boxrule=0pt, arc=1.8mm,
    colback=#1!4!white,
    colbacktitle=#1,
    coltitle=white,
    fonttitle=\sffamily\bfseries,
    left=7pt, right=7pt, top=5pt, bottom=5pt,
    toptitle=3pt, bottomtitle=3pt,
    before skip=10pt, after skip=8pt,
  }
}

% ── Style encadré question type ──────────────────────────────────────
% Carte blanche, bandeau sombre, liseré très clair.
\tcbset{
  questiontype/.style={
    breakable, enhanced jigsaw,
    colback=white,
    frame hidden, boxrule=0pt, arc=1.8mm,
    colbacktitle=colElec!78!black,
    coltitle=white,
    fonttitle=\sffamily\bfseries\small,
    left=7pt, right=7pt, top=4pt, bottom=4pt,
    toptitle=2.5pt, bottomtitle=2.5pt,
    before skip=8pt, after skip=8pt,
    borderline={0.6pt}{0pt}{colElec!12!white},
  }
}

% ── Style encadré jury ───────────────────────────────────────────────
% Bandeau ambre, fond crème, sans cadre.
\tcbset{
  jury/.style={
    breakable, enhanced jigsaw,
    colback=colJuryBg,
    frame hidden, boxrule=0pt, arc=1.8mm,
    colbacktitle=colJuryFrame,
    coltitle=white,
    fonttitle=\sffamily\bfseries,
    left=7pt, right=7pt, top=4pt, bottom=4pt,
    toptitle=3pt, bottomtitle=3pt,
    before skip=10pt, after skip=8pt,
  }
}

% ── En-tête / pied de page communs (refonte 2026) ─────────────────────
% #1 code fiche, #2 titre court, #3 couleur domaine
\newcommand{\ficheHeaders}[3]{%
  \fancyhead[L]{\small\sffamily{\bfseries\textcolor{#3}{#1}}\hspace{0.7em}\textcolor{colSlate}{#2}}%
  \fancyhead[R]{\small\sffamily\textcolor{colSlate}{CAPET SII IE\, --- \,Session 2026}}%
  \fancyfoot[C]{\small\sffamily\textcolor{colSlate}{\thepage\,/\,\pageref{LastPage}}}%
  \fancyhead[C]{}\fancyfoot[L]{}\fancyfoot[R]{}%
}

% ── Bloc en-tête de fiche (BLOC 0, refonte 2026) ──────────────────────
% #1 code  #2 titre  #3 couleur domaine  #4 domaine  #5 ligne présence  #6 sous-thèmes
\newcommand{\ficheEntete}[6]{%
  \begin{tcolorbox}[enhanced, breakable,
    frame hidden, boxrule=0pt, arc=2.5mm,
    interior style={left color=#3, right color=#3!75!black},
    left=9pt, right=9pt, top=7pt, bottom=7pt,
    before skip=0pt, after skip=6pt]
    {\color{white}\sffamily\bfseries\LARGE #1\quad #2}\\[4pt]
    {\color{#3!22!white}\sffamily\small \MakeUppercase{#4}\quad\textbar\quad #5}\\[1pt]
    {\color{#3!35!white}\sffamily\footnotesize Sous-thèmes : #6}
  \end{tcolorbox}%
}

% ── Mise en page ─────────────────────────────────────────────────────
\pagestyle{fancy}
\fancyhf{}
\renewcommand{\headrulewidth}{0.5pt}
\renewcommand{\footrulewidth}{0pt}
\AtBeginDocument{%
  \renewcommand{\headrule}{{\color{colRule}\hrule width\headwidth height 0.5pt}}%
}

\setlist[itemize]{leftmargin=1.2em, itemsep=1pt, topsep=2pt, parsep=0pt}
\setlist[enumerate]{leftmargin=1.2em, itemsep=3pt, topsep=2pt, parsep=0pt}

\setlength{\parskip}{2pt}
\setlength{\columnsep}{10pt}

\def\CAPETBEGINFILE{\begin{document}}
\def\CAPETENDFILE{\end{document}}
\else
\def\CAPETBEGINFILE{}
\def\CAPETENDFILE{}
\fi

\newcommand{\totalPages}{5}  % mis - jour apres premiere compile

\ficheHeaders{T02}{Puissance en régime AC}{colElec}

\CAPETBEGINFILE

% -----------------------------------------------------------------------
%  BLOC EN-TETE (integre - la premiere page, sans saut de page)
% -----------------------------------------------------------------------

\ficheEntete{T02}{Puissance en régime alternatif sinusoïdal (AC)}{colElec}%
  {électrotechnique}%
  {\textbf{Présence :} \textbf{13/13} sujets (2013-2026) - thème \emph{systématiquement} présent}%
  {triangle des puissances - facteur de puissance - triphasé équilibré - compensation réactive - bilan de puissance}

% -----------------------------------------------------------------------
%  1. RAPPELS DE COURS
% -----------------------------------------------------------------------

\begin{tcolorbox}[
  ficheSection=colElec,
  title={1. Rappels de cours},
]

% -- 1.1 Monophase --------------------------------------------------
\textbf{1.1 Régime sinusoïdal monophasé}

\smallskip
Récepteur d'impédance $\underline{Z} = R + jX$, alimenté par $u(t) = U\sqrt{2}\cos(\omega t)$.
Courant : $i(t) = I\sqrt{2}\cos(\omega t - \varphi)$, avec $\varphi = \arg(\underline{Z})$.

\smallskip
\begin{center}
\renewcommand{\arraystretch}{1.45}
\begin{tabular}{llll}
\toprule
\textbf{Grandeur} & \textbf{Symb.} & \textbf{Expression} & \textbf{Unité} \\
\midrule
Puissance active    & $P$   & $P = UI\cos\varphi = RI^2 = \dfrac{U^2\cos\varphi}{Z}$ & W \\[5pt]
Puissance réactive  & $Q$   & $Q = UI\sin\varphi = XI^2 = \dfrac{U^2\sin\varphi}{Z}$ & VAR \\[5pt]
Puissance apparente & $S$   & $S = UI = \sqrt{P^2+Q^2}$ & VA \\[4pt]
Facteur de puissance & $\cos\varphi$ & $\cos\varphi = \dfrac{P}{S} = \dfrac{R}{Z}$ & - $\in[0;1]$ \\[5pt]
Puissance complexe  & $\underline{S}$ & $\underline{S} = \underline{U}\,\overline{\underline{I}} = P + jQ$ & VA \\
\bottomrule
\end{tabular}
\end{center}

\smallskip
\textbf{Triangle des puissances :}

\begin{center}
\begin{minipage}[c]{0.32\textwidth}
% Schema TikZ
\begin{tikzpicture}[scale=1.05, font=\small]
  \draw[thick, colElec, ->] (0,0) -- (3.2,0) node[right] {$P$};
  \draw[thick, colElec!60!black, ->] (3.2,0) -- (3.2,1.8) node[right] {$Q$};
  \draw[thick, colElec!40!black, ->] (0,0) -- (3.2,1.8) node[above left=-2pt] {$S$};
  % angle phi
  \draw[colElec!70!black] (0.7,0) arc[start angle=0, end angle=29.3, radius=0.7];
  \node at (1.05, 0.2) {\footnotesize$\varphi$};
  % labels
  \node[below] at (1.6,0) {$P = S\cos\varphi$};
  \node[right] at (3.25,0.9) {$Q = S\sin\varphi$};
\end{tikzpicture}
\end{minipage}
\hfill
\begin{minipage}[c]{0.30\textwidth}
\[S^2 = P^2 + Q^2\]
\[\tan\varphi = \frac{Q}{P}\]
\end{minipage}
\hfill
\begin{minipage}[c]{0.30\textwidth}
\footnotesize
$Q > 0$ : inductif (moteur)\\
$Q < 0$ : capacitif (condo)\\
$Q = 0$ : resistif, $\cos\varphi=1$
\end{minipage}
\end{center}

\medskip
\hrule\vspace{4pt}

% -- 1.2 Triphase --------------------------------------------------
\textbf{1.2 Régime triphasé équilibré}

$V$ = tension simple (phase-neutre), $U$ = tension composée (entre phases) :
$\displaystyle U = \sqrt{3}\,V$

\smallskip
\begin{center}
\renewcommand{\arraystretch}{1.4}
\begin{tabular}{lll}
\toprule
\textbf{Grandeur} & \textbf{Expression} & \textbf{Remarque} \\
\midrule
Puissance active    & $P = 3VI\cos\varphi = \sqrt{3}\,UI\cos\varphi$ & $I$ = courant de ligne \\[4pt]
Puissance réactive  & $Q = 3VI\sin\varphi = \sqrt{3}\,UI\sin\varphi$ & \\[4pt]
Puissance apparente & $S = 3VI = \sqrt{3}\,UI = \sqrt{P^2+Q^2}$ & \\[4pt]
Facteur de puissance & $\cos\varphi = P/S$ & identique au monophasé \\
\bottomrule
\end{tabular}
\end{center}

\smallskip
\textbf{Couplage étoile (Y) / triangle ($\Delta$) :}

\smallskip
\begin{center}
\renewcommand{\arraystretch}{1.25}
\begin{tabular}{lcc}
\toprule
 & \textbf{étoile (Y)} & \textbf{Triangle ($\Delta$)} \\
\midrule
Tension aux bornes du recepteur & $V = U/\sqrt{3}$ & $U$ \\
Courant dans le recepteur & $I_\text{ligne}$ & $I_\text{ligne}/\sqrt{3}$ \\
\bottomrule
\end{tabular}
\end{center}

\medskip
\hrule\vspace{4pt}

% -- 1.3 Compensation ----------------------------------------------
\textbf{1.3 Compensation de l'énergie réactive}

Objectif : rôlever $\cos\varphi_1$ - $\cos\varphi_2 > \cos\varphi_1$ en ajoutant des condensateurs en parallèle.

\[
  Q_C = P\,\bigl(\tan\varphi_1 - \tan\varphi_2\bigr)
  \qquad \text{[puissance réactive - fournir par les condensateurs]}
\]

\begin{center}
\renewcommand{\arraystretch}{1.3}
\begin{tabular}{lll}
\toprule
\textbf{Couplage} & \textbf{Tension sur le condo} & \textbf{Capacite par phase} \\
\midrule
étoile (Y) & $V = U/\sqrt{3}$ & $\displaystyle C = \dfrac{Q_C}{3\,\omega\,V^2}$ \\[6pt]
Triangle ($\Delta$) & $U$ & $\displaystyle C = \dfrac{Q_C}{3\,\omega\,U^2}$ \\[4pt]
\bottomrule
\end{tabular}
\end{center}
\vspace{2pt}
{\footnotesize Sur réseau 400\,V/230\,V : $V=230$\,V, $\omega = 2\pi\times50 \approx 314$\,rad/s. Un condensateur $\Delta$ vaut 3 fois moins qu'un equivalent Y.}

\medskip
\hrule\vspace{4pt}

% -- 1.4 Bilan installation ----------------------------------------
\textbf{1.4 Bilan de puissance d'une installation complexe}

Plusieurs recepteurs en parallèle - \textbf{sommer $P$ et $Q$ séparément, jamais $S$} :
\[
  P_\text{tot} = \textstyle\sum P_k \qquad Q_\text{tot} = \textstyle\sum Q_k \qquad S_\text{tot} = \sqrt{P_\text{tot}^2 + Q_\text{tot}^2}
\]
Pour un moteur asynchrone ($P_u$ = puissance utile mecanique, $\eta$ = rendement) :
\[
  P_a = \frac{P_u}{\eta} \qquad
  Q = P_a\tan\varphi \qquad
  S = \frac{P_a}{\cos\varphi}
\]

\end{tcolorbox}

% -----------------------------------------------------------------------
%  2. FORMULAIRE SYNTHeTIQUE (2 colonnes)
% -----------------------------------------------------------------------

\begin{tcolorbox}[
  ficheSection=colElec,
  title={2. Formulaire synthétique},
]

\small
\begin{multicols}{2}
\renewcommand{\arraystretch}{1.45}
\begin{tabular}{@{}p{3.8cm}ll@{}}
\toprule
\textbf{Formule} & \textbf{Usage} & \textbf{U.} \\
\midrule
\multicolumn{3}{@{}l}{\footnotesize\textit{Monophase}} \\
$P = UI\cos\varphi$ & P active & W \\
$Q = UI\sin\varphi$ & P réactive & VAR \\
$S = \sqrt{P^2+Q^2}$ & P apparente & VA \\
$\cos\varphi = P/S = R/Z$ & Facteur puis. & - \\
\midrule
\multicolumn{3}{@{}l}{\footnotesize\textit{Triphase équilibré}} \\
$P = \sqrt{3}\,UI\cos\varphi$ & P active & W \\
$Q = \sqrt{3}\,UI\sin\varphi$ & P réactive & VAR \\
$S = \sqrt{3}\,UI$ & P apparente & VA \\
$U = \sqrt{3}\,V$ & Tension & V \\
$I = P/(\sqrt{3}\,U\cos\varphi)$ & Courant ligne & A \\
\bottomrule
\end{tabular}

\columnbreak

\renewcommand{\arraystretch}{1.45}
\begin{tabular}{@{}p{3.8cm}ll@{}}
\toprule
\textbf{Formule} & \textbf{Usage} & \textbf{U.} \\
\midrule
\multicolumn{3}{@{}l}{\footnotesize\textit{Compensation réactive}} \\
$Q_C = P(\tan\varphi_1 - \tan\varphi_2)$ & - compenser & VAR \\[2pt]
$C = Q_C/(3\omega V^2)$ & Condo étoile & F \\[2pt]
$C = Q_C/(3\omega U^2)$ & Condo triangle & F \\
\midrule
\multicolumn{3}{@{}l}{\footnotesize\textit{Moteur asynchrone}} \\
$P_a = P_u / \eta$ & P absorbée & W \\
$Q = P_a \tan\varphi$ & P réactive & VAR \\
$S = P_a / \cos\varphi$ & P apparente & VA \\
\midrule
\multicolumn{3}{@{}l}{\footnotesize\textit{Trigonometrie}} \\
$\sin\varphi = \sqrt{1-\cos^2\!\varphi}$ & & - \\
$\tan\varphi = \sin\varphi/\cos\varphi$ & & - \\
\bottomrule
\end{tabular}

\end{multicols}
\normalsize

\end{tcolorbox}

% -----------------------------------------------------------------------
%  3. QUESTIONS TYPES
% -----------------------------------------------------------------------

\begin{tcolorbox}[
  ficheSection=colElec,
  title={3. Questions types - extraites des sujets réels},
]

% -- QT1 -------------------------------------------------------------

\begin{tcolorbox}[questiontype,
  title={QT1 - Bilan de puissance d'une installation triphasée \hfill\normalfont\footnotesize E2 disciplinaire 2024}
]

\textbf{énoncé.}
Un atelier triphasé $U=\SI{400}{\volt}$, $f=\SI{50}{\hertz}$ comporte :
un moto-extracteur ($I_n=\SI{229}{\ampere}$ relevé sur la plaque signalétique, $\cos\varphi_\text{mot}=0{,}85$) ;
des ordinateurs ($P_0=\SI{1}{\kilo\watt}$, $\cos\varphi_0=1$) ;
eclairage et équipements ($P_e=\SI{43}{\kilo\watt}$, $Q_e=\SI{25{,}5}{\kilo VAR}$).
\textbf{Calculer} $P_\text{ch}$, $Q_\text{ch}$, $S_\text{ch}$, $I_\text{abs}$ et $F_p$.

\tcbline

\textbf{Corrigé.}

\textbf{étape 1 - puissances du moteur} (à partir des données de plaque $I_n$, $\cos\varphi_\text{mot}$).
\[
  P_\text{mot} = \sqrt{3}\,U I_n \cos\varphi_\text{mot} = \sqrt{3}\times 400\times 229\times 0{,}85 \approx \SI{134{,}9}{\kilo\watt}
\]
\[
  \sin\varphi_\text{mot} = \sqrt{1-0{,}85^2} = 0{,}527
  \qquad
  Q_\text{mot} = \sqrt{3}\,U I_n \sin\varphi_\text{mot} \approx \SI{83{,}6}{\kilo VAR}
\]

\textbf{étape 2 - bilan global} (sommer $P$ et $Q$ séparément, jamais $S$).
\[
  P_\text{ch} = 134{,}9 + 1 + 43 = \SI{178{,}9}{\kilo\watt}
  \qquad
  Q_\text{ch} = 83{,}6 + 0 + 25{,}5 = \SI{109{,}1}{\kilo VAR}
\]
Le corrigé officiel impose pour la suite $P_\text{ch}=\SI{178{,}3}{\kilo\watt}$ et
$Q_\text{ch}=\SI{108{,}7}{\kilo VAR}$ (écart d'arrondi sur les données de plaque) --- valeurs
retenues ci-après.
\[
  S_\text{ch} = \sqrt{P_\text{ch}^2+Q_\text{ch}^2} = \sqrt{178{,}3^2+108{,}7^2} \approx \SI{208{,}8}{\kilo VA}
\]

\textbf{étape 3 - courant et facteur de puissance.}
\[
  I_\text{abs} = \frac{S_\text{ch}}{\sqrt{3}\,U} = \frac{208\,800}{\sqrt{3}\times 400} \approx \SI{301}{\ampere}
  \qquad
  F_p = \frac{P_\text{ch}}{S_\text{ch}} = \frac{178{,}3}{208{,}8} \approx \mathbf{0{,}854}
\]

\textbf{Conclusion.} $F_p \approx 0{,}85 < 0{,}92$ : l'installation ne respecte pas l'exigence du
gestionnaire d'énergie ($F_p \in [0{,}92\,;\,1]$) --- une compensation de l'énergie réactive est
nécessaire (cf.\ QT2).

\end{tcolorbox}

\vspace{4pt}

% -- QT2 -------------------------------------------------------------

\begin{tcolorbox}[questiontype,
  title={QT2 - Compensation de l'énergie réactive par condensateurs \hfill\normalfont\footnotesize E2 disciplinaire 2024}
]

\textbf{énoncé.}
(Données jury, identiques au bilan de QT1 : $P=\SI{178{,}3}{\kilo\watt}$, $Q=\SI{108{,}7}{\kilo VAR}$.)
Un gestionnaire impose $F_p=1$ via une batterie de condensateurs étoile ($C$ par phase).
\textbf{B3.} Montrer que $C = Q_\text{abs}/(3\omega V_n^2)$.
\textbf{B4.} Calculer $C$.

\tcbline

\textbf{Corrigé.}

\textbf{B3.}
Pour $F_p'=1$ : $Q_\text{total} = Q_\text{abs} + Q_C = 0$, donc $Q_C = -Q_\text{abs}$.
Un condensateur étoile sous tension simple $V_n$ fournit $Q_C = 3C\omega V_n^2$, d'où :
\[
  \boxed{C = \frac{Q_\text{abs}}{3\,\omega\,V_n^2}}
\]

\textbf{B4.}
$V_n = \SI{230}{\volt}$, $\omega = 2\pi\times 50 \approx \SI{314}{\radian\per\second}$ :
\[
  C = \frac{108\,700}{3 \times 314 \times 230^2} = \frac{108\,700}{49\,831\,800} \approx \SI{2{,}18}{\milli\farad}
\]

\textbf{Conclusion.} Avec $C \approx \SI{2{,}2}{\milli\farad}$ par phase en étoile, la puissance
réactive absorbée est intégralement compensée : $F_p = 1$, conforme à l'exigence du gestionnaire
d'énergie ($F_p \in [0{,}92\,;\,1]$).

\end{tcolorbox}

\vspace{4pt}

% -- QT3 -------------------------------------------------------------

\begin{tcolorbox}[questiontype,
  title={QT3 - Relevement du facteur de puissance vers une valeur cible \hfill\normalfont\footnotesize E2 2017-2020}
]

\textbf{énoncé.}
Installation triphasée : $P=\SI{120}{\kilo\watt}$, $\cos\varphi_1=0{,}75$.
On souhaite $\cos\varphi_2=0{,}93$.
Calculer $Q_C$ et $C$ (couplage étoile, $V=\SI{230}{\volt}$).

\tcbline

\textbf{Corrigé.}
\[
  \tan\varphi_1 = \frac{\sqrt{1-0{,}75^2}}{0{,}75} = 0{,}882
  \qquad
  \tan\varphi_2 = \frac{\sqrt{1-0{,}93^2}}{0{,}93} = 0{,}396
\]
\[
  Q_C = P(\tan\varphi_1-\tan\varphi_2) = 120\times(0{,}882-0{,}396) = \SI{58{,}3}{\kilo VAR}
\]
\[
  C = \frac{58\,300}{3\times 314\times 230^2} \approx \SI{1{,}17}{\milli\farad}
\]
\textit{Vérif. :}
\[
  \cos\varphi_2 = \frac{120}{\sqrt{120^2+(105{,}8-58{,}3)^2}} = \frac{120}{129{,}0} \approx 0{,}930 \;\checkmark
\]

\textbf{Conclusion.} Une batterie étoile de $C \approx \SI{1{,}17}{\milli\farad}$ par phase fournit
les \SI{58{,}3}{\kilo VAR} requis : la vérification redonne bien $\cos\varphi_2 \approx 0{,}93$,
l'objectif de relèvement du facteur de puissance fixé par le cahier des charges est atteint.

\end{tcolorbox}

\end{tcolorbox}

% -----------------------------------------------------------------------
%  4. ATTENTES DU JURY
% -----------------------------------------------------------------------

\begin{tcolorbox}[
  jury,
  title={4. Attentes du jury - Erreurs récurrentes et points de vigilance},
]

\begin{enumerate}[label=\textbf{\arabic*.}, leftmargin=*, itemsep=5pt]

  \item \textbf{Ne jamais sommer les puissances apparentes $S_k$.}
  $S_\text{tot} \neq S_1+S_2+\cdots$ \textit{(erreur eliminatoire)}.
  Méthode : sommer $P_k$, sommer $Q_k$, puis $S=\sqrt{P^2+Q^2}$.

  \item \textbf{Confondre tension simple et composée.}
  $V = U/\sqrt{3}$ - réseau 400\,V/230\,V : $U=400$\,V entre phases, $V=230$\,V phase-neutre.
  Les formules $P=\sqrt{3}\,UI\cos\varphi$ utilisent $U$ composée et $I$ de ligne.

  \item \textbf{Oublier le rendement du moteur.}
  $P_\text{absorbée} = P_\text{utile}/\eta$. C'est $P_a$ qui sert - calculer $Q$ et $S$, pas $P_u$.

  \item \textbf{Erreurs d'unités et d'ordre de grandeur.}
  Le jury signale chaque année des erreurs de facteur $10^3$.
  Convertir systématiquement en SI avant toute application numerique.

  \item \textbf{Confondre couplage étoile et triangle pour $C$.}
  étoile : $C = Q_C/(3\omega V^2)$ avec $V$ simple.
  Triangle : $C = Q_C/(3\omega U^2)$ avec $U$ composée.
  Un condensateur $\Delta$ vaut 3 fois moins qu'un condensateur Y equivalent.

  \item \textbf{Signe de la puissance réactive.}
  Inductances ($Q>0$), condensateurs ($Q<0$).
  La compensation annule la somme algebrique des $Q_k$.

  \item \textbf{Rédiger la demonstration, pas seulement le resultat.}
  Exemple attendu : \textit{"Pour $F_p=1$, $Q_\text{tot}=Q_\text{abs}+Q_C=0$, donc $Q_C=-Q_\text{abs}$.
  Un condensateur étoile donne $Q_C=3C\omega V_n^2$, d'où $C=Q_\text{abs}/(3\omega V_n^2)$."}

  \item \textbf{Verifier le resultat.}
  Contrôler que $P\leq S$, $\cos\varphi\in[0;1]$, valeur de $C$ realiste (mF - eF selon les puissances).

\end{enumerate}

\medskip
\begin{tcolorbox}[
  colback=colJuryFrame!12!colJury,
  colframe=colJuryFrame,
  left=4pt, right=4pt, top=2pt, bottom=2pt,
  fontupper=\small\itshape
]
Rapports jury 2019-2025 : "La manipulation des unités et des puissances de 10 est tres souvent hasardeuse, entraînant des resultats dont l'ordre de grandeur meme est faux."
"Un nombre significatif de candidats ne parvient pas - caracteriser des constituants electriques ou - identifier leur rôle dans une chaîne de puissance."
\end{tcolorbox}

\end{tcolorbox}

% -----------------------------------------------------------------------
%  PIED - PRESENCE DANS LES SUJETS
% -----------------------------------------------------------------------

\begin{tcolorbox}[
  ficheSection=colPresence,
  title={Présence dans les sujets analyses (2013-2026)},
]
\begin{center}
\footnotesize\renewcommand{\arraystretch}{1.15}
\begin{tabular}{lccccccccccccc}
\toprule
\textbf{Annee} & 13 & 14 & 15 & 16 & 17 & 18 & 19 & 20 & 21 & 22 & 23 & 24 & 25 \\
\midrule
E1 (commune)    & $\checkmark$ & $\checkmark$ & $\checkmark$ & $\checkmark$ & $\checkmark$ & $\checkmark$ & $\checkmark$ & $\checkmark$ & $\checkmark$ & $\checkmark$ & $\checkmark$ & $\checkmark$ & $\checkmark$ \\
E2 discipl. IE  & - & - & - & - & $\checkmark$ & $\checkmark$ & $\checkmark$ & $\checkmark$ & $\checkmark$ & $\checkmark$ & $\checkmark$ & $\checkmark$ & $\checkmark$ \\
\bottomrule
\end{tabular}
\end{center}
\smallskip
\footnotesize\textbf{Sous-thèmes frequents :}
$P/Q/S$ triphasé (13/13) - facteur de puissance (13/13) - bilan installation (11/13) - compensation condensateurs (9/13) - chute de tension ligne (6/13)

\medskip
\textbf{Ordres de grandeur (repères rapides).}
\begin{itemize}[leftmargin=1.3em]
  \item Reseau BT FR : tension phase-neutre typique \SI{230}{\volt}, tension phase-phase typique \SI{400}{\volt}, fréquence \SI{50}{\hertz}.
  \item Courants d'etude frequents en sujet : de quelques \si{\ampere} a quelques dizaines de \si{\ampere}.
  \item En conversion d'énergie, un rendement attendu reste generalement entre 0{,}8 et 0{,}98 selon la technologie.
\end{itemize}

\medskip
\textbf{Synthese express (à refaire sans support).}
\begin{enumerate}[leftmargin=1.4em]
  \item Poser les hypotheses du modele et les conventions de signe.
  \item Ecrire 2 a 4 relations clés avec unités.
  \item Réaliser un calcul numerique complet et verifier l'ordre de grandeur.
  \item Justifier le resultat par un argument technique de niveau CAPET.
\end{enumerate}

\end{tcolorbox}

\CAPETENDFILE










